\documentclass[11pt]{article}
\usepackage[margin=1in]{geometry}
\usepackage{setspace}
\usepackage{hyperref}
\hypersetup{colorlinks=true, allcolors=black}
\setlength{\parindent}{0pt}
\setlength{\parskip}{0.55em}
\setstretch{1.05}
\pagestyle{empty}

\begin{document}

\begin{center}
{\Large Teaching Statement}\\[0.35em]
Ke Cheng HSU (Clio)\\[0.15em]
{\small Department of Economics and Finance, City University of Hong Kong}
\end{center}

I have taught at the Hong Kong University of Science and Technology since 2020, at both undergraduate and postgraduate level. From 2020 to 2025 I was a teaching assistant for undergraduate macroeconomics, the economics of uncertainty and information, and Ph.D.\ microeconomic theory. I led tutorials, held office hours, graded, and for some courses designed the tutorial materials and wrote examination questions. My average evaluation across these appointments was 4.72 out of 5. In 2026 I designed and lectured the mathematics camp for the MSc in Economics, covering the calculus, linear algebra, optimization, and dynamic-programming background that students need for econometrics and theory. Many students came from non-mathematics backgrounds, so I built each topic step by step and used a picture whenever it could carry the idea, such as the geometry of projection for ordinary least squares. The instructor evaluation was 4.53 out of 5. As Teaching Assistant Coordinator in the Business School in 2025--2026, I helped postgraduate students design teaching materials and evaluated their mock teaching.

I am ready to teach economics at undergraduate and postgraduate level, including evening classes, and to supervise student projects. The courses closest to my training are empirical economics, applied econometrics, and quantitative courses that use data from the Chinese economy. My own research estimates how return policy and pricing work on a large Taobao apparel panel, and how retail pricing and unsold inventory behave in Japanese bookstore data. In class I would use that kind of evidence to show how a regression, a panel, or a simple structural model answers a policy or business question, and only then write down the formula. I can also teach principles and intermediate microeconomics and macroeconomics, mathematics for economists, and an introduction to empirical research practice: cleaning data, keeping a project under version control, and checking results rather than trusting generated code.

My classroom rule is that students should be able to say, in plain language, why an outcome occurs. A long queue at a new campus restaurant is not proof that the food is good. It is a signal, and later students can rationally join even when their own information says otherwise. That example opens herding and information cascades before any formula. I prepare study materials and problem sets that follow the same order, and I write examination questions that ask students to apply the mechanism, not only to recall a definition. I ask them to stop me when a step is unclear, because one question usually belongs to several people. Teaching is most useful when a student who was stuck can explain the result back, and then use it on a project of their own.

\end{document}
